% from ega4-16-fr.tex
\oldpage[IV-4]{40}
pour $n\leq m$ (16.7.7) qu'un opérateur différentiel d'ordre $\leq n$ est
aussi un opérateur différentiel d'ordre $\leq m$ pour tout $m\geq n$. Si
$D:\mathscr{F}\to\mathscr{G}$ est un opérateur différentiel d'ordre
$\leq n$, alors, pour tout ouvert $U$ de $X$,
$D|U:\mathscr{F}|U\to\mathscr{G}|U$ est aussi un opérateur différentiel
d'ordre $\leq n$.

On dit qu'un homomorphisme $D:\mathscr{F}\to\mathscr{G}$ des faisceaux de
groupes additifs sous-jacents à $\mathscr{F}$ et à $\mathscr{G}$ est un
\emph{opérateur différentiel} (relativement à $S$) si, pour tout $x\in X$, il
existe un voisinage ouvert $U$ de $x$ et un entier $n\geq0$ tels que
$D|U:\mathscr{F}|U\to\mathscr{G}|U$ soit un opérateur différentiel d'ordre
$\leq n$. L'\emph{ordre} d'un opérateur différentiel
$D:\mathscr{F}\to\mathscr{G}$ est la borne inférieure des entiers $n$ tels
que $D$ soit un opérateur différentiel d'ordre $\leq n$ (et donc $+\infty$
s'il n'y a pas de tels entiers); cet ordre est toujours fini si $X$ est
\emph{quasi-compact}. Les opérateurs différentiels d'ordre $0$ ne sont autres
que les homomorphismes de $\mathcal{O}_X$-Modules
$\mathscr{F}\to\mathscr{G}$; on convient que tout opérateur différentiel
d'ordre $<0$ est \emph{nul}. Pour $n\geq0$, un opérateur différentiel n'est
pas en général un homomorphisme de $\mathcal{O}_X$-Modules mais est toujours
un homomorphisme de $\psi^*(\mathcal{O}_S)$-Modules.

Lorsque $\mathscr{F}=\mathcal{O}_X$, un opérateur différentiel d'ordre
$\leq1$ de $\mathcal{O}_X$ dans $\mathscr{G}$ s'écrit d'une seule manière
sous la forme $v+D$, où $v:\mathcal{O}_X\to\mathscr{G}$ est un
$\mathcal{O}_X$-homomorphisme, et $D$ une $S$-\emph{dérivation} (16.5.1) de
$\mathcal{O}_X$ dans $\mathscr{G}$: cela résulte de la structure de
$P_{B/A}^1$ (0, 20.4.8).

\begin{env}[16.8.2]
\label{IV.16.8.2}
Pour décrire de façon plus explicite un opérateur différentiel d'ordre
$\leq n$, $D:\mathscr{F}\to\mathscr{G}$, il suffit, pour tout ouvert affine
$U$ de $X$, dont l'image dans $S$ est contenue dans un ouvert affine $V$, de
caractériser l'homomorphisme
$D=D_U:\Gamma(U,\mathscr{F})\to\Gamma(U,\mathscr{G})$. Si l'on pose
$\Gamma(V,\mathcal{O}_S)=A$, $\Gamma(U,\mathcal{O}_X)=B$, de sorte que $B$
est une $A$-algèbre, on a
\[
  \Gamma(U,\mathscr{P}_{X/S}^n)
  =(B\otimes_A B)/\mathfrak{I}^{n+1},
\]
où l'on a posé pour abréger $\mathfrak{I}=\mathfrak{I}_{B/A}$. Posons en
outre $M=\Gamma(U,\mathscr{F})$, $N=\Gamma(U,\mathscr{G})$; alors la
définition de $D$ signifie que pour chaque couple $(U,V)$ vérifiant les
conditions précédentes, le $A$-homomorphisme $D:M\to N$ se factorise en
\[
  M\longrightarrow
  ((B\otimes_A B)/\mathfrak{I}^{n+1})\otimes_B M
  \xrightarrow{v}N
\]
où la première flèche est l'homomorphisme canonique $t\mapsto1\otimes t$, et
$v$ est un $B$-homomorphisme, la structure de $B$-module de
$((B\otimes_A B)/\mathfrak{I}^{n+1})\otimes_B M$ provenant du premier
facteur $B$ (alors qu'on rappelle que dans la formation du produit tensoriel
sur $B$, la structure de $B$-module de
$(B\otimes_A B)/\mathfrak{I}^{n+1}$ provient du second facteur $B$). Notons
maintenant que le $B$-module
$((B\otimes_A B)/\mathfrak{I}^{n+1})\otimes_B M$ est isomorphe à
$(B\otimes_A M)/\mathfrak{I}^{n+1}(B\otimes_A M)$, où $B\otimes_A M$ est
considéré comme $(B\otimes_A B)$-module et sa structure de $B$-module
provient de l'homomorphisme $b\mapsto b\otimes1$ de $B$ dans $B\otimes_A B$.
Soit alors $D'$ le $B$-homomorphisme de $B\otimes_A M$ dans $N$ tel que
$D'(b\otimes t)=bD(t)$; la condition de factorisation sur $D$ s'exprime encore
en disant que $D'$ doit \emph{être nul dans le $B$-module}
$\mathfrak{I}^{n+1}(B\otimes_A M)$.
\end{env}

\begin{env}[16.8.3]
\label{IV.16.8.3}
Il est clair que l'ensemble des opérateurs différentiels d'ordre $\leq n$ de
$\mathscr{F}$ dans $\mathscr{G}$ est un groupe additif, noté
$\operatorname{Diff}_{X/S}^n(\mathscr{F},\mathscr{G})$; lorsque
$\mathscr{F}=\mathscr{G}=\mathcal{O}_X$, on écrit aussi
$\operatorname{Diff}_{X/S}^n$ au lieu de
$\operatorname{Diff}_{X/S}^n(\mathcal{O}_X,\mathcal{O}_X)$.

On a vu (16.8.1) que l'on a pour deux ouverts $U\supset V$ de $X$, un
homomorphisme canonique de restriction
\[
  \operatorname{Diff}_{U/S}^n(\mathscr{F}|U,\mathscr{G}|U)
  \longrightarrow
  \operatorname{Diff}_{V/S}^n(\mathscr{F}|V,\mathscr{G}|V)
\]
